\section{Related Work}
\subsection{Diffusion Models}
Diffusion models constitute a major advancement in generative modeling, facilitating notable progress in image synthesis, video generation, and molecular design~\citep{ho2020denoising, dhariwal2021diffusion, croitoru2022vision}. These models function by gradually adding noise to data and then learning to reverse this process to generate new samples~\citep{song2020score}.
However, the iterative denoising process of diffusion models leads to high computational costs, limiting their applicability to edge devices and real-time scenarios~\citep{zheng2025diffusion,ho2020denoising}. To address this issue, numerous studies have sought to reduce the number of sampling steps by developing more advanced ODE solvers. For instance, Denoising Diffusion Implicit Models (DDIMs) propose a non-Markovian formulation of the diffusion process, thereby substantially reducing the required sampling steps~\citep{song2020denoising}.

% Their strength lies in producing high-quality samples via iterative denoising, but this comes at a high computational cost due to the many sampling steps. To address this, research has explored acceleration techniques such as improved variance schedules, advanced ODE solvers, and non-Markovian formulations like DDIMs~\citep{song2020denoising, nichol2021improved, liu2022pseudo}, which reduce sampling time while largely preserving quality. Nonetheless, balancing efficiency and fidelity remains an open challenge.

\subsection{Modulated Diffusion}
In contrast to approaches that reduce the number of denoising steps, quantization aims to enhance the inference efficiency of diffusion models by encoding floating-point values as low-bit integers~\citep{nagel2021white, yang2019quantization}. Post-training quantization (PTQ) is particularly appealing because it eliminates the need for retraining~\citep{li2021brecq}. However, PTQ faces significant challenges in activation quantization due to the presence of activation outliers and dynamic input ranges~\citep{li2023q, wang2024towards, huang2024tfmq, shang2023post, he2023ptqd, zhao2024vidit, xiao2023smoothquant, feng2024mpq}. Additionally, discrepancies in bitwidth between weights and activations further diminish the computational efficiency of quantization on hardware~\citep{gope2020high}.

% Acceleration techniques for diffusion models include caching, which reuses intermediate computations across time steps~\citep{ma2024deepcache, wimbauer2024cache, ma2024learning}, and quantization, which reduces the precision of weights and activations~\citep{nagel2021white, yang2019quantization}. Post-training quantization (PTQ) is attractive because it avoids retraining~\citep{li2021brecq}, but diffusion models suffer from activation outliers and time-varying ranges, restricting PTQ to 8-bit activations~\citep{li2023q, wang2024towards, huang2024tfmq, shang2023post, he2023ptqd, zhao2024vidit, xiao2023smoothquant, feng2024mpq}. Quantization-aware training (QAT) improves low-bit performance but requires retraining~\citep{nagel2022overcoming, feng2024mpq, he2023efficientdm}.

Modulated Diffusion Models (MoDiff)~\citep{gao2025modulateddiffusionacceleratinggenerative} introduce a general framework to address activation-quantization challenges in PTQ methods. By exploiting redundancy throughout the generation process and incorporating error compensation mechanisms, MoDiff reduces quantization error and prevents error accumulation. As a result, MoDiff reduces activation quantization from 8 bits to 4 bits without compromising performance relative to existing baselines. Nevertheless, the original evaluation of MoDiff is limited to theoretical bit-operation (BOP) analysis and does not include empirical hardware-acceleration results.

% \red{Do we need some references for kernel design?}

% MoDiff introduces modulated quantization, which accelerates diffusion models by exploiting temporal differences in activations~\citep{gao2025modulateddiffusionacceleratinggenerative}. It computes activations incrementally by quantizing differences between consecutive time steps, which have smaller ranges and fewer outliers than raw activations, enabling lower-bit quantization. To avoid error accumulation, MoDiff uses error-compensated modulation that tracks quantization errors and corrects subsequent computations. This reduces quantization error in theory and supports quantization down to 3 bits without performance loss. However, the original work only analyzed simulated FLOPs without implementation or hardware experiments. This work extends MoDiff with a full implementation, real-time performance measurements, and comprehensive evaluations of generation quality and efficiency.